\documentclass[12pt,a4paper]{article}
\PassOptionsToPackage{hidelinks}{hyperref}

% Shared IHEDN report formatting: cover, typography, tables and bibliography.
\usepackage{ihedn-report}

% Example bibliography supplied with the template.
\addbibresource{vendor/biblatex-sciences-po-fr/bibliographies/articles.bib}
\addbibresource{vendor/biblatex-sciences-po-fr/bibliographies/livres.bib}
\addbibresource{vendor/biblatex-sciences-po-fr/bibliographies/online.bib}
\addbibresource{vendor/biblatex-sciences-po-fr/bibliographies/rapports.bib}
\addbibresource{vendor/biblatex-sciences-po-fr/bibliographies/theses.bib}

\usepackage[acronym,toc,section=section]{glossaries}
\makeglossaries
\newacronym{codir}{CODIR}{Comité de direction}
\newacronym{retex}{RETEX}{Retour d expérience}
\newacronym{rgpd}{RGPD}{Règlement général sur la protection des données}
\renewcommand*{\acronymname}{Liste des acronymes}

% Shared data and rendering used by the template cover example.
% Shared cover definition used by main.tex and cover-example.tex.
\newcommand{\IhednTemplateCommittee}{Comité d'étude, de mois AAAA à mois AAAA, composé de : Monsieur / Madame Nom Prénom, \textit{président ou présidente} ; Monsieur / Madame Nom Prénom, \textit{rapporteur ou rapporteure} ; Monsieur / Madame Nom Prénom ; Monsieur / Madame Nom Prénom ; Monsieur / Madame Nom Prénom.}

\newcommand{\IhednTemplateCover}{%
  % Citation optionnelle : décommenter et adapter la ligne suivante si nécessaire.
  % \Citation{Texte de la citation}[Auteur de la citation]
  \President{Monsieur / Madame Nom Prénom}{Président ou Présidente du comité}%
  \Rapporteur{Monsieur / Madame Nom Prénom}{Rapporteur ou Rapporteure du comité}%
  \PageDeGardeIHEDN{
    header_1={Union-IHEDN},
    header_2={Association des auditeurs IHEDN},
    header_3={région Paris / Île de France},
    association_line={AR~16 : Association régionale Paris Île de France},
    committee_label={Comité de rédaction de l'AR~16},
    date={Soumis le : JJ mois AAAA},
    logo_path={Figures/logos_IHEDN/Logo_AR_16_PARIS_IDF.jpg},
    title={Titre du rapport : sous-titre éventuel},
    title_max_lines=2,
    title_fontsize={26.000pt},
    title_baselineskip={28.667pt},
    forum_line_1={Contribution pour le Forum des Études de l'UNION -- IHEDN AAAA},
    forum_line_2={{\itshape\enquote{Esprit de Défense \& Engagement dans la France d'aujourd'hui}}},
    committee={\IhednTemplateCommittee}
  }%
}


\begin{document}

% ----------------------------------------------------------------
% Cover shared with cover-example.tex.
% ----------------------------------------------------------------
\IhednTemplateCover

% Front matter mirrors the structure of a complete IHEDN report.
\pagenumbering{Roman}
\addcontentsline{toc}{section}{Résumé}
\section*{Résumé}
Ce document fournit un exemple complet des éléments disponibles pour rédiger un rapport IHEDN : page de garde, propositions, sommaire, acronymes, figures, tableaux, annexes et bibliographie.

\listofpropositions
\clearpage
\tableofcontents
\clearpage
\printglossary[type=\acronymtype,title={Liste des acronymes}]
\clearpage
\addcontentsline{toc}{section}{Table des figures}
\clearpage
\listoffigures
\clearpage
\addcontentsline{toc}{section}{Liste des tableaux}
\begingroup
\makeatletter
\renewcommand{\@pnumwidth}{3em}
\makeatother
\listoftables
\endgroup
\clearpage
\pagenumbering{arabic}

\section{La rédaction du rapport}

\begin{itemize}
  \item Le mémoire fera entre 20 et 40~pages, en police Times New Roman, caractère~12, auxquelles se rajoutent les annexes.
  \item Le rapport comprend un sommaire et s'organise en chapitres ou sections, qui traitent chacun ou chacune un aspect de la démonstration générale.
  \item Le compte rendu doit démontrer une thèse ou \textit{a minima} apporter une plus-value~: il est donc nécessaire de construire un plan démonstratif, qui permette d'argumenter en direction de l'hypothèse formulée.
  \item Défendre une thèse suppose d'en formuler clairement les termes. Le raisonnement doit ensuite être appuyé sur des éléments empiriques aussi précis, consistants et soigneusement établis que possible.
  \item Il ne faut pas chercher à masquer les lacunes des matériaux empiriques à disposition~: il faut compenser ces faiblesses en croisant les différentes sources pour réunir un faisceau d'indices, mais en soulignant, lorsque c'est nécessaire le statut encore hypothétique de certaines des conclusions émises. Dans le cadre de ce travail exploratoire, c'est tout à fait acceptable.
  \item À chaque chapitre ou section, il est nécessaire de présenter des propositions concrètes et réalisables.
\end{itemize}

Le corps d'un rapport peut employer des acronymes tels que le \acrshort{codir}, le \gls{retex} ou le \gls{rgpd}. Une référence peut être présentée directement en note de bas de page\footcite{EC:2022:AAR}.

\begin{proposition}{Formuler une recommandation claire et réalisable}
Chaque proposition doit être reliée au diagnostic et préciser les acteurs concernés, les moyens nécessaires et les résultats attendus.
\end{proposition}

\propositiongroup{Décliner une proposition structurante}
\begin{subproposition}{Préciser la première modalité de mise en œuvre}
Présenter une action vérifiable, son responsable et son calendrier.
\end{subproposition}
\begin{subproposition}{Prévoir le suivi et l'évaluation}
Définir les critères qui permettront d'apprécier les effets de la proposition.
\end{subproposition}


\newpage

\section{Insertion de figures}

Chaque figure doit comporter une légende achevée par un point et être dimensionnée de manière à respecter strictement les marges du document.
Chaque figure doit être référencée dans le texte au moyen de son numéro~: \ref{fig:exemple}.

\begin{figure}[h]
    \centering
    \includegraphics[width=0.8\textwidth]{example-image.pdf}
    \caption{Titre explicite de la figure.}
    \label{fig:exemple}
\end{figure}



\begin{figure}[htbp]
    \centering
    \begin{subfigure}[b]{0.48\textwidth}
        \centering
        \includegraphics[width=\textwidth]{example-image-a.pdf}
        \caption{Première illustration.}
        \label{fig:exemple-a}
    \end{subfigure}
    \hfill
    \begin{subfigure}[b]{0.48\textwidth}
        \centering
        \includegraphics[width=\textwidth]{example-image-b.pdf}
        \caption{Seconde illustration.}
        \label{fig:exemple-b}
    \end{subfigure}
    \caption{Exemple de figure composée de deux sous-figures.}
    \label{fig:exemple-compose}
\end{figure}

\begin{table}[H]
    \centering
    \StyledTableSetup
    \begin{tabularx}{\linewidth}{@{}p{0.28\linewidth}Y Y@{}}
        \toprule
        \textbf{Étape} & \textbf{Objectif} & \textbf{Livrable} \\
        \midrule
        Cadrage & Définir le périmètre et la méthode. & Note de cadrage validée. \\
        Analyse & Croiser les données et les entretiens. & Diagnostic argumenté. \\
        Proposition & Formuler une réponse réalisable. & Plan d'action et indicateurs. \\
        \bottomrule
    \end{tabularx}
    \caption{Exemple de tableau structuré.}
    \label{tab:exemple-structure}
\end{table}



\section{Diagrammes et planification}

\begin{figure}[H]
    \centering
    \begin{tikzpicture}[
        node distance=18mm,
        block/.style={draw,rounded corners,align=center,minimum width=38mm,minimum height=10mm}
    ]
        \node[block] (diagnostic) {Diagnostic};
        \node[block,right=of diagnostic] (decision) {Décision};
        \node[block,right=of decision] (action) {Mise en œuvre};
        \draw[-{Stealth}] (diagnostic) -- (decision);
        \draw[-{Stealth}] (decision) -- (action);
    \end{tikzpicture}
    \caption{Exemple de chaîne de décision réalisée avec TikZ.}
    \label{fig:exemple-tikz}
\end{figure}

\begin{figure}[H]
    \centering
    \begin{ganttchart}[
        hgrid,
        vgrid,
        x unit=1.45cm,
        y unit chart=0.65cm,
        bar height=0.55
    ]{1}{6}
        \gantttitle{Calendrier prévisionnel}{6} \\
        \gantttitlelist{1,...,6}{1} \\
        \ganttbar{Cadrage}{1}{2} \\
        \ganttbar{Analyse}{2}{4} \\
        \ganttbar{Restitution}{5}{6}
    \end{ganttchart}
    \caption{Exemple de planification réalisée avec PGF Gantt.}
    \label{fig:exemple-gantt}
\end{figure}

\begin{landscape}
\section{Tableau long en orientation paysage}
\StyledTableSetup
\begin{longtable}{@{}p{0.18\linewidth}p{0.34\linewidth}p{0.34\linewidth}@{}}
    \caption{Exemple de tableau long pouvant se poursuivre sur plusieurs pages.}
    \label{tab:exemple-long} \\
    \toprule
    \textbf{Phase} & \textbf{Travaux} & \textbf{Résultats attendus} \\
    \midrule
    \endfirsthead
    \toprule
    \textbf{Phase} & \textbf{Travaux} & \textbf{Résultats attendus} \\
    \midrule
    \endhead
    Préparation & Définition du périmètre, des objectifs et des responsabilités. & Gouvernance et calendrier validés. \\
    Collecte & Réunion des données documentaires et conduite des entretiens. & Corpus vérifié et traçable. \\
    Analyse & Croisement des sources, formulation du diagnostic et discussion des limites. & Constats argumentés et hiérarchisés. \\
    Proposition & Élaboration des recommandations et de leurs modalités de suivi. & Plan d'action exploitable. \\
    \bottomrule
\end{longtable}
\end{landscape}

\newpage

\section{Références et bibliographie}

La bibliographie doit suivre les points détaillés ci-dessous.

\begin{itemize}

  \item Le rapport doit obligatoirement comprendre des références et une bibliographie. \textbf{Il faut référencer avec la plus grande rigueur tous les travaux utilisés.}
  \item Dans le cadre de ce travail, le référencement en note de bas de page en indiquant la référence \textit{in extenso} garantit une lisibilité optimale pour les relecteurs et une aisance de rédaction accrue.
  \item Il faut impérativement établir une bibliographie complète des sources scientifiques en fin de volume en suivant le même mode de présentation.
  \item Il est possible en cas de bibliographie importante (supérieure à 6~pages) de proposer un classement bibliographique («~études généralistes~», «~articles de presse~», «~articles universitaires~», «~ouvrages~», etc.).
\end{itemize}

\subsection{Usage des locutions latines}

Certaines locutions latines, traditionnellement utilisées dans les références
bibliographiques, peuvent être utilisées afin d'abréger les citations successives
et alléger la lecture.

\begin{table}[h]
	\centering
	\renewcommand{\arraystretch}{1.3}
	\begin{tabular}{|>{\centering\arraybackslash}p{3cm}|>{\centering\arraybackslash}p{6.5cm}|}
		\hline
		\textbf{Locution} &
		\textbf{Condition d'apparition} \\
		\hline

		\emph{Ibid.} &
		Deux citations consécutives renvoyant à la \textbf{même œuvre}. 
		La pagination peut être identique ou modifiée. \\
		\hline

		\emph{Id.} / \emph{Ead.} &
		Deux citations consécutives renvoyant au \textbf{même auteur,}
		mais à des œuvres différentes. \\
		\hline

		\emph{Op. cit.} &
		Retour à une œuvre déjà citée après \textbf{interruption}
		par une autre référence. \\
		\hline

		\emph{Loc. cit.} &
		Retour à la \textbf{même œuvre et au même passage}
		dans un contexte non consécutif. \\
		\hline
	\end{tabular}
	\caption{Usage des principales locutions latines dans les références bibliographiques.}
  \label{table:locutions_latines}
\end{table}

\newpage

\subsection{Livres (imprimés ou \textit{e-book})}

\fullcite{Benjamin:2018}

\fullcite{Castanet:2023}

\fullcite{Pilkington:1998}

\subsection{Rapports (imprimés ou en ligne)}

\fullcite{EC:2022:AAR}

\fullcite{Claeys:2005:PolicyMix}

\subsection{Chapitres d'ouvrage}

\fullcite{Blancheton:2017:Crises:Online}

\subsection{Articles de revue}

\fullcite{Rinde:Rokkan:1959:Online}

\subsection{Articles de presse (en ligne)}

\fullcite{Lechner:2014:Eotone}

\subsection{Thèses et mémoires (imprimé ou numérique)}

\fullcite{Delanoue:2018}

\subsection{Statistiques}

\fullcite{GlobalWebIndex:2020}

\subsection{Vidéos en ligne}

\fullcite{Has:2016:Zotero}

\subsection{Podcasts}

\fullcite{FranceCulture:2023:Holodomor}

\subsection{Pages web}

\fullcite{PresidenceRepublique:2018:BoutiqueElysee}

\clearpage
\appendix
\section{Exemple d'annexe}
Les annexes accueillent les documents complémentaires nécessaires à la compréhension du rapport sans alourdir le corps de la démonstration.

\subsection{Inclusion d'un document PDF}
L'utilitaire partagé peut intégrer toutes les pages d'un PDF dans des cadres homogènes :

\IhednRenderPdfPagesInBoxes{reference/page-de-garde.pdf}

\clearpage
\addcontentsline{toc}{section}{Références bibliographiques}
\printbibheading[title={Références bibliographiques}]

\addcontentsline{toc}{subsection}{Ouvrages}
\newrefcontext[sorting=nyt]
\printbibliography[type=book,title={Ouvrages},heading=subbibliography]
\printbibliography[type=incollection,title={Chapitres d'ouvrage},heading=subbibliography]

\addcontentsline{toc}{subsection}{Articles de recherche}
\newrefcontext[sorting=nyt]
\printbibliography[type=article,title={Articles de recherche},heading=subbibliography]

\addcontentsline{toc}{subsection}{Rapports institutionnels}
\newrefcontext[sorting=nyt]
\printbibliography[type=report,title={Rapports institutionnels},heading=subbibliography]

\addcontentsline{toc}{subsection}{Ressources en ligne}
\newrefcontext[sorting=nyt]
\printbibliography[type=online,title={Ressources en ligne},heading=subbibliography]

\addcontentsline{toc}{subsection}{Thèses et mémoires}
\newrefcontext[sorting=nyt]
\printbibliography[type=thesis,title={Thèses et mémoires},heading=subbibliography]

% Compilation: LATEXMK_SHELL_ESCAPE=0 latexmk -pdf main.tex

\end{document}
